% GENERATED FILE. Edit canonical lessons, then run npm run generate:books.
% canonical-source-hash: fnv1a64:eb2a7e9f4b7f67dd
% canonical-lessons: PA-C295-tarer, PA-C295-phatna, PA-C295-khokhla, PA-C295-khadd, PA-C295-galna

\chapter{Cracks and Hollows}
\label{ch:pa-a2-cracks-and-hollows}
\hlchaptermodality{295}

\begin{chapteropening}
\textbf{By the end of this chapter:} \emph{I can describe cracked or hollow things.}

Complete the last lesson of chapter 295: I can describe cracked or hollow things.
\end{chapteropening}

\section[\texorpdfstring{\pa{ਤਰੇੜ} (tareṛ)}{ਤਰੇੜ (tareṛ)}]{\pa{ਤਰੇੜ} (tareṛ) — a crack}

\label{lesson:PA-C295-tarer}



\begin{quote}
Before the new one: say the Punjabi for a clod, then the Punjabi for dust.
\end{quote}

\subsection*{You'll want to know: \pa{ਤਰੇੜ}}
\textbf{\pa{ਤਰੇੜ}} — \emph{tareṛ} — \textquotedblleft{}a crack\textquotedblright{}.

\begin{quote}
\pa{ਕੰਧ} \pa{ਵਿੱਚ} \pa{ਇੱਕ} \pa{ਲੰਮੀ} \pa{ਤਰੇੜ} \pa{ਹੈ।}
\end{quote}

\begin{quote}
\emph{kandh vich ikk lammī tareṛ hai.}
\end{quote}

\begin{quote}
There is a long crack in the wall.
\end{quote}

A feminine noun.

\begin{cousinweb}
\textbf{\pa{ਤਰੇੜ}} (\emph{tareṛ}) is inherited from a reconstructed \emph{*tṛdi}, \textquotedblleft{}crack, fissure\textquotedblright{}; CDIAL lists Punjabi \emph{tareṛ}, \textquotedblleft{}crack, crevice\textquotedblright{}, under it (CDIAL 5926).
\end{cousinweb}

\begin{grammarlens}[title={What you've built}]
One more word for this chapter.
\end{grammarlens}

\subsection*{Your turn}
\begin{itemize}
\raggedright
  \item \emph{Say it:} \emph{tareṛ}
  \item \emph{Say it:} \emph{tareṛ}, once more
  \item \emph{Say it:} \emph{kandh vich ikk lammī tareṛ hai.} — There is a long crack in the wall.
  \item \emph{Recall:} say the Punjabi for a clod, then the Punjabi for dust, then say \emph{tareṛ} again
\end{itemize}

\begin{tcolorbox}[breakable,colback=teal!4,colframe=teal!35!black,title={Before you move on}]
What does \pa{ਤਰੇੜ} mean? (\textquotedblleft{}A crack\textquotedblright{}.) Say it once more.
\end{tcolorbox}
\section[\texorpdfstring{\pa{ਫਟਣਾ} (phaṭṇā)}{ਫਟਣਾ (phaṭṇā)}]{\pa{ਫਟਣਾ} (phaṭṇā) — to split open; to crack}

\label{lesson:PA-C295-phatna}



\begin{quote}
Before the new one: say the Punjabi for dust, then the Punjabi for a crack.
\end{quote}

\subsection*{You'll want to know: \pa{ਫਟਣਾ}}
\textbf{\pa{ਫਟਣਾ}} — \emph{phaṭṇā} — \textquotedblleft{}to split open; to crack\textquotedblright{}.

\begin{quote}
\pa{ਸੋਕੇ} \pa{ਵਿੱਚ} \pa{ਜ਼ਮੀਨ} \pa{ਫਟਦੀ} \pa{ਹੈ।}
\end{quote}

\begin{quote}
\emph{soke vich zamīn phaṭdī hai.}
\end{quote}

\begin{quote}
In a drought the ground splits open.
\end{quote}

Past (it …): \pa{ਉਹ} \pa{ਘੜਾ} \pa{ਫਟਿਆ}, \emph{oh ghaṛā phaṭiā}.

\begin{cousinweb}
\textbf{\pa{ਫਟਣਾ}} (\emph{phaṭṇā}) goes back to Sanskrit \emph{sphāṭayati}, \textquotedblleft{}causes to split\textquotedblright{}, as a word taken from Hindi: CDIAL lists Punjabi \emph{phāṭṇā}, \textquotedblleft{}to be split\textquotedblright{}, under it, and derives the short-vowel \emph{phaṭṇā} from Hindi \emph{phaṭnā} (CDIAL 13825).
\end{cousinweb}

\begin{grammarlens}[title={What you've built}]
One more word for this chapter.
\end{grammarlens}

\subsection*{Your turn}
\begin{itemize}
\raggedright
  \item \emph{Say it:} \emph{phaṭṇā}
  \item \emph{Say it:} \emph{phaṭṇā}, once more
  \item \emph{Say it:} \emph{soke vich zamīn phaṭdī hai.} — In a drought the ground splits open.
  \item \emph{Recall:} say the Punjabi for dust, then the Punjabi for a crack, then say \emph{phaṭṇā} again
  \item \emph{Write it:} copy \textbf{\pa{ਫਟਣਾ}} once, checking the whole word
\end{itemize}

\begin{tcolorbox}[breakable,colback=teal!4,colframe=teal!35!black,title={Before you move on}]
What does \pa{ਫਟਣਾ} mean? (\textquotedblleft{}To split open; to crack\textquotedblright{}.) Say it once more.
\end{tcolorbox}
\section[\texorpdfstring{\pa{ਖੋਖਲਾ} (khokhlā)}{ਖੋਖਲਾ (khokhlā)}]{\pa{ਖੋਖਲਾ} (khokhlā) — hollow}

\label{lesson:PA-C295-khokhla}



\begin{quote}
Before the new one: say the Punjabi for a crack, then the Punjabi for to split open.
\end{quote}

\subsection*{You'll want to know: \pa{ਖੋਖਲਾ}}
\textbf{\pa{ਖੋਖਲਾ}} — \emph{khokhlā} — \textquotedblleft{}hollow\textquotedblright{}.

\begin{quote}
\pa{ਇਹ} \pa{ਪੁਰਾਣਾ} \pa{ਰੁੱਖ} \pa{ਖੋਖਲਾ} \pa{ਹੈ।}
\end{quote}

\begin{quote}
\emph{eh purāṇā rukkh khokhlā hai.}
\end{quote}

\begin{quote}
This old tree is hollow.
\end{quote}

\begin{cousinweb}
\textbf{\pa{ਖੋਖਲਾ}} (\emph{khokhlā}) is inherited from a reconstructed \emph{*khōkkha}, \textquotedblleft{}hollow\textquotedblright{}; CDIAL lists Punjabi \emph{khokhlā} under it (CDIAL 3927).
\end{cousinweb}

\begin{grammarlens}[title={What you've built}]
One more word for this chapter.
\end{grammarlens}

\subsection*{Your turn}
\begin{itemize}
\raggedright
  \item \emph{Say it:} \emph{khokhlā}
  \item \emph{Say it:} \emph{khokhlā}, once more
  \item \emph{Say it:} \emph{eh purāṇā rukkh khokhlā hai.} — This old tree is hollow.
  \item \emph{Recall:} say the Punjabi for a crack, then the Punjabi for to split open, then say \emph{khokhlā} again
  \item \emph{Write it:} copy \textbf{\pa{ਖੋਖਲਾ}} once, checking the whole word
\end{itemize}

\begin{tcolorbox}[breakable,colback=teal!4,colframe=teal!35!black,title={Before you move on}]
What does \pa{ਖੋਖਲਾ} mean? (\textquotedblleft{}Hollow\textquotedblright{}.) Say it once more.
\end{tcolorbox}
\section[\texorpdfstring{\pa{ਖੱਡ} (khaḍḍ)}{ਖੱਡ (khaḍḍ)}]{\pa{ਖੱਡ} (khaḍḍ) — a ravine; a deep pit}

\label{lesson:PA-C295-khadd}



\begin{quote}
Before the new one: say the Punjabi for to split open, then the Punjabi for hollow.
\end{quote}

\subsection*{You'll want to know: \pa{ਖੱਡ}}
\textbf{\pa{ਖੱਡ}} — \emph{khaḍḍ} — \textquotedblleft{}a ravine; a deep pit\textquotedblright{}.

\begin{quote}
\pa{ਪਹਾੜ} \pa{ਦੇ} \pa{ਹੇਠਾਂ} \pa{ਡੂੰਘੀ} \pa{ਖੱਡ} \pa{ਹੈ।}
\end{quote}

\begin{quote}
\emph{pahāṛ de heṭhāṁ ḍūṅghī khaḍḍ hai.}
\end{quote}

\begin{quote}
There is a deep ravine below the mountain.
\end{quote}

A feminine noun.

\begin{cousinweb}
\textbf{\pa{ਖੱਡ}} (\emph{khaḍḍ}) is inherited from a reconstructed \emph{*khaḍḍa}, \textquotedblleft{}hole, pit\textquotedblright{}; CDIAL lists Punjabi \emph{khaḍḍ}, \textquotedblleft{}pit, ravine\textquotedblright{}, under it (CDIAL 3790).
\end{cousinweb}

\begin{grammarlens}[title={What you've built}]
One more word for this chapter.
\end{grammarlens}

\subsection*{Your turn}
\begin{itemize}
\raggedright
  \item \emph{Say it:} \emph{khaḍḍ}
  \item \emph{Say it:} \emph{khaḍḍ}, once more
  \item \emph{Say it:} \emph{pahāṛ de heṭhāṁ ḍūṅghī khaḍḍ hai.} — There is a deep ravine below the mountain.
  \item \emph{Recall:} say the Punjabi for to split open, then the Punjabi for hollow, then say \emph{khaḍḍ} again
  \item \emph{Write it:} copy \textbf{\pa{ਖੱਡ}} once, checking the whole word
\end{itemize}

\begin{tcolorbox}[breakable,colback=teal!4,colframe=teal!35!black,title={Before you move on}]
What does \pa{ਖੱਡ} mean? (\textquotedblleft{}A ravine; a deep pit\textquotedblright{}.) Say it once more.
\end{tcolorbox}
\section[\texorpdfstring{\pa{ਗਲਣਾ} (galṇā)}{ਗਲਣਾ (galṇā)}]{\pa{ਗਲਣਾ} (galṇā) — to rot; to decay}

\label{lesson:PA-C295-galna}



\begin{quote}
Before the new one: say the Punjabi for hollow, then the Punjabi for a ravine.
\end{quote}

\subsection*{You'll want to know: \pa{ਗਲਣਾ}}
\textbf{\pa{ਗਲਣਾ}} — \emph{galṇā} — \textquotedblleft{}to rot; to decay\textquotedblright{}.

\begin{quote}
\pa{ਗਿੱਲੇ} \pa{ਫਲ} \pa{ਜਲਦੀ} \pa{ਗਲਦੇ} \pa{ਹਨ।}
\end{quote}

\begin{quote}
\emph{gille phal jaldī galde han.}
\end{quote}

\begin{quote}
Wet fruit rots quickly.
\end{quote}

Past (it …): \pa{ਉਹ} \pa{ਸੇਬ} \pa{ਗਲਿਆ}, \emph{oh seb galiā}.

\begin{cousinweb}
\textbf{\pa{ਗਲਣਾ}} (\emph{galṇā}) is inherited from Sanskrit \emph{galati}, \textquotedblleft{}drips, oozes\textquotedblright{}; CDIAL lists Punjabi \emph{galṇā}, \textquotedblleft{}to melt, sink\textquotedblright{}, under it; the sense \textquotedblleft{}to rot\textquotedblright{} is not among those it gives (CDIAL 4074).
\end{cousinweb}

\begin{grammarlens}[title={What you've built}]
One more word for this chapter.
\end{grammarlens}

\subsection*{Your turn}
\begin{itemize}
\raggedright
  \item \emph{Say it:} \emph{galṇā}
  \item \emph{Say it:} \emph{galṇā}, once more
  \item \emph{Say it:} \emph{gille phal jaldī galde han.} — Wet fruit rots quickly.
  \item \emph{Recall:} say the Punjabi for hollow, then the Punjabi for a ravine, then say \emph{galṇā} again
  \item \emph{Write it:} copy \textbf{\pa{ਗਲਣਾ}} once, checking the whole word
\end{itemize}

\begin{tcolorbox}[breakable,colback=teal!4,colframe=teal!35!black,title={Before you move on}]
What does \pa{ਗਲਣਾ} mean? (\textquotedblleft{}To rot; to decay\textquotedblright{}.) Say it once more.
\end{tcolorbox}
